% !TeX root = ../Book4.tex
\begin{figure}[htbp]
	\centering
	\begin{tikzpicture}[font=\small,x=1.05cm,y=1.05cm]
		\fill[GeometryBlueLight] (-3,0) rectangle (3,1.65);
		\fill[orange!12] (-3,-1.65) rectangle (3,0);
		\draw[->,GeometryGray] (-3.3,0)--(3.4,0) node[right] {\(x\)};
		\draw[->,GeometryGray] (0,-1.9)--(0,1.95) node[above] {\(y\)};
		\foreach \x in {-2.7,-2.1,-1.5,-0.9,-0.3,0.3,0.9,1.5,2.1,2.7}
			\fill[GeometryBlue] (\x,0) circle (1.7pt);
		\node[fill=white,inner sep=4pt] at (0,1.03) {\(y>0\)：一张二维叶};
		\node[fill=white,inner sep=4pt] at (0,-1.03) {\(y<0\)：一张二维叶};
		\node[fill=white,inner sep=3pt,anchor=west] at (3.55,0) {\shortstack{轴上每一点\\各是一张零维叶}};
	\end{tikzpicture}
	\caption[秩跳跃的平面 Poisson 结构]{\(\pi=y\partial_x\wedge\partial_y\) 的辛叶。两种底色各表示一个完整的二维叶；横轴上的圆点只列举其中若干零维叶，整条横轴不是一张叶。}
	\label{g4:fig:0904:rank-jump}
\end{figure}
